\mysection{二重否定}

非ゼロ値を1（または論理値\IT{true}）に、ゼロ値を0（または論理値\IT{false}）に変換する一般的な方法\footnote{この方法は読みにくいコードを生成するため議論の余地がある}は\IT{!!statement}です：

\begin{lstlisting}[style=customc]
int convert_to_bool(int a)
{
	return !!a;
};
\end{lstlisting}

最適化GCC 5.4 x86：

\begin{lstlisting}[style=customasmx86]
convert_to_bool:
        mov     edx, DWORD PTR [esp+4]
        xor     eax, eax
        test    edx, edx
        setne   al
        ret
\end{lstlisting}

\INS{XOR}は、\INS{SETNE}がトリガーされない場合でも、常に\EAX内の戻り値をクリアします。
つまり、\INS{XOR}はデフォルトの戻り値を0に設定します。

\myindex{x86!\Instructions!SET}
入力値がゼロと等しくない場合（\INS{SET}命令の\TT{-NE}サフィックス）、
1が\ALに設定され、そうでなければ\ALは変更されません。

なぜ\INS{SETNE}は\EAXレジスタの下位8ビット部分で動作するのでしょうか？
なぜなら重要なのは最後のビット（0または1）だけで、他のビットは\INS{XOR}によってクリアされているからです。

したがって、そのC/C++コードは次のように書き換えることができます：

\begin{lstlisting}[style=customc]
int convert_to_bool(int a)
{
	if (a!=0)
		return 1;
	else
		return 0;
};
\end{lstlisting}

\dots または：

\begin{lstlisting}[style=customc]
int convert_to_bool(int a)
{
	if (a)
		return 1;
	else
		return 0;
};
\end{lstlisting}

\INS{SET}に類似した命令を持たない\ac{CPU}をターゲットとするコンパイラは、この場合、
分岐命令などを生成します。
